\documentclass[10pt,a4paper]{amsart}
\usepackage[margin=19mm]{geometry}
\usepackage{amsmath,amssymb,mathtools}
\usepackage[scr=rsfs,cal=euler]{mathalfa}
\usepackage{xfrac}
\usepackage[unicode,hidelinks,bookmarks=false]{hyperref}
\newcommand{\ie}{i.e.\ }
\newcommand{\cf}{cf.\ }
\newcommand{\resp}{respectively\ }
\newcommand{\Subsection}{\subsection}
\providecommand{\nobreakdash}{\nobreak\hbox{-}}
\setlength{\emergencystretch}{2em}
\multlinegap0pt
\usepackage{amssymb}
\usepackage{mathtools}
\usepackage[shortlabels]{enumitem}
\usepackage{xcolor}
\usepackage{tikz}
\usepackage{float}
\newtheorem{prop}{Proposition}
\title{Remarks on Topological Rigidity of Real Moment-Angle Manifolds}
\author{Ioannis Gkeneralis}
\date{Source: arXiv:2604.15462v1 (CC BY 4.0)}
\begin{document}
\maketitle

\noindent\emph{Source and licence.} Remarks on Topological Rigidity of Real Moment-Angle Manifolds. Version arXiv:2604.15462v1 (CC BY 4.0), distributed by arXiv under the Creative Commons Attribution 4.0 International licence, http://creativecommons.org/licenses/by/4.0/, as recorded in the arXiv licence metadata for this exact version and shown on the abstract page https://arxiv.org/abs/2604.15462v1. Full licence notice: \texttt{licenses/arxiv-2604.15462-CC-BY-4.0.txt}.\par\smallskip

\noindent\textit{Editorial scope.} Counted unit: the article's prop prop:cat0\_action (source0.tex lines 545--552). The article's complete original proof of it is delivered with it (source0.tex lines 553--563, one \texttt{proof} environment of 11 lines). No mathematics is written, repaired or re-derived: the statement and the proof are the article's own lines, in the article's own notation, and no retained line is substituted or re-flowed.  Retained uncounted as the source-local context the proof needs: lines 519--526.  Nothing was omitted from any retained span, so the package carries no omission record.  No retained line contains a citation, so the wrapper carries no bibliography and no bibliography entry.  No retained line contains a figure environment, an image-inclusion command or drawing code, so no image is delivered and the package declares no asset dependency.  Editorial formatting only, in addition: an AMS \texttt{amsart} preamble that restates the article's own declarations and macros used by the retained lines, namely the environments \texttt{prop} on separate counters, together with the standard packages those lines need.  The retained environments print in this document as Proposition 1 in order of appearance, whereas the article prints the same items with the numbering of its own full document.  The complete native source of the article as distributed in the arXiv e-print for this exact version is bundled unchanged as \texttt{sources/198607/source0.tex}.
\begin{prop}
\label{prop:universal-cover-davis}
Let $K$ be a finite simplicial complex. Then the universal cover of $\mathcal{R}_K$ is $W_K$-equivariantly homeomorphic to the basic construction
\[
U(W_K,K(K)),
\]
where $K(K)$ denotes the canonical cubical chamber associated to $K$.
\end{prop}

\begin{prop}
\label{prop:cat0_action}
Let $K$ be a finite flag simplicial complex such that $\mathcal{R}_K$ is a closed manifold. Then the group
\[
G=\pi_1(\mathcal{R}_K)
\]
acts properly, cocompactly, and isometrically on a finite-dimensional CAT(0) space.
\end{prop}

\begin{proof}
By Proposition~\ref{prop:universal-cover-davis}, the universal cover $\widetilde{\mathcal{R}_K}$ is identified with the Davis complex $\Sigma_K$.

In the right-angled case, $\Sigma_K$ admits a natural piecewise Euclidean cubical metric in which each cube is isometric to a standard Euclidean cube. The link of each vertex is isomorphic to the simplicial complex $K$.

If $K$ is flag, then every vertex link is a flag simplicial complex. Hence, by Gromov's link condition, $\Sigma_K$ is nonpositively curved. Since $\Sigma_K$ is simply connected, it follows that it is a CAT(0) space.

The action of $G=\pi_1(\mathcal{R}_K)$ on $\widetilde{\mathcal{R}_K} \cong \Sigma_K$ is the deck transformation action. Because $\mathcal{R}_K$ is compact, this action is properly discontinuous and cocompact. Moreover, the cubical metric is invariant under the Coxeter group action and hence under its subgroup $G$, so the action is by isometries.

Finally, $\Sigma_K$ is finite-dimensional since $K$ is finite.
\end{proof}

\end{document}
